\subsection*{The Next Construction}

\begin{remark*}[Exposition]
The passage from $\mathbb{N}$ to $\mathbb{W}$ adds the element needed for
$0+a=a$. It does not add solutions to equations such as $a+x=b$ when
$b<a$. Thus $\mathbb{W}$ is still not closed under subtraction.
\end{remark*}

\begin{remark*}[Exposition]
The construction of $\mathbb{Z}$ repairs this by adjoining additive inverses
formally. One standard route uses ordered pairs from
$\mathbb{W}\times\mathbb{W}$, with a pair $(a,b)$ interpreted as a formal
difference $a-b$. The equivalence relation then identifies different pairs
that represent the same intended difference.
\end{remark*}

\begin{dependencies}
\begin{itemize}
  \item \hyperref[thm:w-additive-commutative-monoid]{$(\mathbb{W},+,0)$ Is a Commutative Monoid}
  \item \hyperref[thm:w-has-no-additive-inverses-except-zero]{$\mathbb{W}$ Has No Additive Inverses Except $0$}
  \item \hyperref[thm:w-is-not-a-ring]{$\mathbb{W}$ Is Not a Ring}
\end{itemize}
\end{dependencies}
